\section{A. 경기대학교 AI 컴퓨터공학부 \& 솔루티오}

\begin{frame} % No title at first slide
    \sectiontitle{A}{경기대학교 AI 컴퓨터공학부 \& 솔루티오}
    \sectionmeta{
        출제자 -- 전민주({\href{https://solved.ac/profile/mingmingmon}{@mingmingmon}})\\
        태그: \consolas {\#math, \#arithmetic}\\
        예상 난이도 -- \textbf{\color{acB5}Bronze5}
    }
    \begin{itemize}
        \item 제출 32번, 정답 25명 (정답률 78.1\%)
        \item 처음 푼 사람: 김0건, 김0민, 김0원, 박0민, 안0현 1분
    \end{itemize}
\end{frame}
%----------------------------------------------------------------------------
\begin{frame}{\textbf{A}. 경기대학교 AI 컴퓨터공학부 \& 솔루티오}
    \begin{itemize}
        \item $n$이 19,471,108 이하의 자연수이기 때문에 \\
        나머지 연산자 모듈러(\%)를 사용하여 짝수인지, 홀수인지 판단하면 됩니다.
        \item 2로 나눈 나머지가 0이라면 짝수,\\
        나머지가 1이라면 홀수이므로\\
        각각 "KGU SOLUTIO" 또는 "KGU AI CSE"를 출력하면 됩니다.
        \item if문을 사용하셔도 되고, 한줄로 표현하고 싶다면\\
        삼항 연산자로 표현하셔도 좋습니다.
    \end{itemize}
\end{frame}